
Our methodology follows a controlled experimental design to isolate and measure the relationship between data quality metrics Q(D) and information density.

\subsubsection{Data Quality Metric Q(D)}

We define data quality as a composite of four components, each capturing a distinct aspect of information content:

\textbf{Deduplication Ratio} measures redundancy removal. Duplicate or near-duplicate documents carry zero marginal information. We compute the deduplication ratio as:

$$
\text{dedup\_ratio} = 1 - \frac{\text{tokens in duplicates}}{\text{total tokens}}
$$

where duplicates are identified via MinHash LSH with Jaccard similarity threshold 0.9. A corpus with 40\% duplicate content has dedup\_ratio = 0.6. Higher values indicate less redundancy.

\textbf{Domain Diversity} measures concept coverage breadth. A corpus drawn from only one domain has narrower information content than one spanning multiple domains. We compute diversity via Shannon entropy over domain distributions:

$$
\text{diversity} = -\sum_{i=1}^{N} p_i \log p_i
$$

where $p_i$ is the fraction of tokens from domain $i$. We use 8 C4 source domains. Uniform distribution maximizes diversity.

\textbf{Perplexity Score} measures signal-to-noise ratio. High-perplexity text is harder for a reference language model to predict, indicating either noise or unusual information. We score documents using GPT-2 small (125M parameters) and normalize:

$$
\text{perplexity\_score} = 1 - \frac{\log(\text{ppl}) - \log(\text{ppl}_{\min})}{\log(\text{ppl}_{\max}) - \log(\text{ppl}_{\min})}
$$

where lower perplexity (more predictable) scores higher.

\textbf{Token Efficiency} measures signal density within documents. Documents with excessive whitespace or boilerplate carry less information per token. We compute:

$$
\text{efficiency} = \frac{\text{content tokens}}{\text{total tokens}}
$$

where content tokens exclude punctuation-only sequences, whitespace runs, and repetitive boilerplate patterns.

\textbf{Composite Q(D):} We combine components via weighted average:

$$
Q(D) = w_1 \cdot \text{dedup\_ratio} + w_2 \cdot \text{diversity} + w_3 \cdot \text{perplexity\_score} + w_4 \cdot \text{efficiency}
$$

where weights $w_i$ are learned via linear regression against information density.

\subsubsection{Information Density Measurement}

Information density quantifies how much learning content each token provides. We measure it through three complementary proxies:

\textbf{Token Entropy} measures average information content per token via Shannon entropy over the unigram distribution:

$$
H_{\text{token}} = -\sum_{t \in V} \frac{c(t)}{N} \log\_2 \frac{c(t)}{N}
$$

Higher entropy indicates more uniform token distribution, implying richer vocabulary usage. We normalize by $log_{2}|V|$ to get density $\in$ [0,1].

\textbf{N-gram Redundancy} measures repetitive structure via 3-gram and 5-gram overlap statistics:

$$
\text{redundancy}_n = \frac{\text{repeated n-grams}}{\text{total n-grams}}
$$

and invert to get diversity\_n = 1 - redundancy\_n. High redundancy indicates formulaic, repetitive text.

\textbf{Semantic Diversity} measures concept coverage via embedding-based analysis. We encode sentences using SentenceTransformer, compute pairwise cosine similarities, and measure average distance:

$$
\text{semantic\_diversity} = 1 - \frac{1}{N(N-1)} \sum_{i \neq j} \text{sim}(e\_i, e\_j)
$$

\textbf{Combined Density:} We average the three proxies (after normalization to [0,1]):

$$
\text{info\_density} = \frac{1}{3}(H_{\text{norm}} + \text{diversity}_{\text{ngram}} + \text{semantic\_diversity})
$$

\subsubsection{Controlled Experimental Design}

To isolate each quality dimension's effect, we create 12 C4 subsets (10GB each, 120GB total) with systematic quality variations:

\begin{itemize}
\item 3 deduplication levels: None (baseline, ~40\% duplicates), Medium (~20\% duplicates), High (<5\% duplicates)
\item 3 diversity levels: Low (single domain), Medium (4 domains), High (8 domains)
\item 3 perplexity levels: Low (bottom 33\% GPT-2 perplexity), Medium (middle 33\%), High (top 33\%)
\item 3 efficiency levels: Low (<50th percentile), Medium (50-75th), High (>75th)
\end{itemize}

We use a fractional factorial design to cover all dimensions without requiring 81 subsets. We sample 12 combinations that maximize orthogonality between factors.

Each subset is created by:
\begin{enumerate}
\item Sampling from C4 (allenai/c4, en split) to select 10GB matching the target quality profile
\item Computing Q(D) metrics via MinHash LSH deduplication detection, domain classification, GPT-2 perplexity scoring, and efficiency analysis
\item Computing information density via token entropy, n-gram redundancy, and semantic diversity
\item Correlation analysis: Pearson and Spearman correlations with p-value testing at $\alpha$ = 0.01
\end{enumerate}

\textbf{Success criterion:} All four Q(D) components achieve r > 0.5 and p < 0.01.

\subsubsection{Reproducibility Testing}

We independently resample three additional C4 subsets at the ''Medium'' quality level and compute coefficient of variation (CV = $\sigma /\mu )$ across the four instances.

\textbf{Success criterion:} CV < 10\% for all components.

\subsubsection{Implementation Details}

All experiments use:
\begin{itemize}
\item Dataset: C4 (allenai/c4, en split)
\item Deduplication: MinHash LSH with 128 hash functions, Jaccard threshold 0.9
\item Perplexity model: GPT-2 small (125M params)
\item Embedding model: SentenceTransformer all-MiniLM-L6-v2
\item Computation: V100 GPUs (32GB), 9.2 GPU-hours total
\end{itemize}
